\documentclass[10pt,a4paper]{amsart}
\usepackage{amsmath,amssymb,amsthm,graphicx}
\usepackage{verbatim}
\newenvironment{reference}{\comment}{\endcomment}
%\newenvironment{reference}{}{}
\newenvironment{slogan}{\comment}{\endcomment}
\newenvironment{history}{\comment}{\endcomment}

% For commutative diagrams we use Xy-pic
\usepackage[all]{xy}

% We use 2cell for 2-commutative diagrams.
\xyoption{2cell}
% Theorem environments.
%
\theoremstyle{plain}
\newtheorem{theorem}[subsection]{Theorem}
\newtheorem{proposition}[subsection]{Proposition}
\newtheorem{lemma}[subsection]{Lemma}

\theoremstyle{definition}
\newtheorem{definition}[subsection]{Definition}
\newtheorem{example}[subsection]{Example}
\newtheorem{exercise}[subsection]{Exercise}
\newtheorem{situation}[subsection]{Situation}

\theoremstyle{remark}
\newtheorem{remark}[subsection]{Remark}
\newtheorem{remarks}[subsection]{Remarks}

\numberwithin{equation}{subsection}

% Macros
%
\def\lim{\mathop{\mathrm{lim}}\nolimits}
\def\colim{\mathop{\mathrm{colim}}\nolimits}
\def\Spec{\mathop{\mathrm{Spec}}}
\def\Hom{\mathop{\mathrm{Hom}}\nolimits}
\def\Ext{\mathop{\mathrm{Ext}}\nolimits}
\def\SheafHom{\mathop{\mathcal{H}\!\mathit{om}}\nolimits}
\def\SheafExt{\mathop{\mathcal{E}\!\mathit{xt}}\nolimits}
\def\Sch{\mathit{Sch}}
\def\Mor{\mathop{\mathrm{Mor}}\nolimits}
\def\Ob{\mathop{\mathrm{Ob}}\nolimits}
\def\Sh{\mathop{\mathit{Sh}}\nolimits}
\def\NL{\mathop{N\!L}\nolimits}
\def\CH{\mathop{\mathrm{CH}}\nolimits}
\def\proetale{{pro\text{-}\acute{e}tale}}
\def\etale{{\acute{e}tale}}
\def\QCoh{\mathit{QCoh}}
\def\Ker{\mathop{\mathrm{Ker}}}
\def\Im{\mathop{\mathrm{Im}}}
\def\Coker{\mathop{\mathrm{Coker}}}
\def\Coim{\mathop{\mathrm{Coim}}}

% Boxtimes
%
\DeclareMathSymbol{\boxtimes}{\mathbin}{AMSa}{"02}

%
% Macros for moduli stacks/spaces
%
\def\QCohstack{\mathcal{QC}\!\mathit{oh}}
\def\Cohstack{\mathcal{C}\!\mathit{oh}}
\def\Spacesstack{\mathcal{S}\!\mathit{paces}}
\def\Quotfunctor{\mathrm{Quot}}
\def\Hilbfunctor{\mathrm{Hilb}}
\def\Curvesstack{\mathcal{C}\!\mathit{urves}}
\def\Polarizedstack{\mathcal{P}\!\mathit{olarized}}
\def\Complexesstack{\mathcal{C}\!\mathit{omplexes}}
% \Pic is the operator that assigns to X its picard group, usage \Pic(X)
% \Picardstack_{X/B} denotes the Picard stack of X over B
% \Picardfunctor_{X/B} denotes the Picard functor of X over B
\def\Pic{\mathop{\mathrm{Pic}}\nolimits}
\def\Picardstack{\mathcal{P}\!\mathit{ic}}
\def\Picardfunctor{\mathrm{Pic}}
\def\Deformationcategory{\mathcal{D}\!\mathit{ef}}
\usepackage{cleveref}
\crefname{originalreference}{original reference}{original references}
\Crefname{originalreference}{Original reference}{Original references}
\makeatletter
\newlabel{lemma-dim-0}{{0FHQ}{source}{}{sourcepointer-1}{}}
\newlabel{lemma-dim-0@cref}{{[originalreference][0][]0FHQ}{[1][0][]source}{}{}{}}
\makeatother
\crefname{originalreference}{original reference}{original references}
\Crefname{originalreference}{Original reference}{Original references}
\makeatletter
\newlabel{lemma-weil-additive}{{0FHJ}{source}{}{sourcepointer-2}{}}
\newlabel{lemma-weil-additive@cref}{{[originalreference][0][]0FHJ}{[1][0][]source}{}{}{}}
\makeatother
\title{Base change of cohomology along a finite separable algebra}
\author{The Stacks Project Authors}
\date{stacks-project:weil.tex:local-snapshot-20261012; GNU FDL 1.2 or later}
\begin{document}
\maketitle
\noindent\textit{Editorial note.} All delivered mathematical text is the author's own, copied from the pinned Stacks Project chapter file weil.tex; no statement or proof step was rewritten, reordered or removed.\par\medskip
\noindent\textit{Editorial note.} The counted proof is the author's complete argument: the decomposition of the tensor product of the finite separable algebras, and the comparison of the cohomologies of the base change through additivity and the product formula.\par\medskip
\noindent\textit{Editorial note.} The label lemma-otimes names the counted statement, and the proof environment immediately following it in the source is delivered verbatim as its proof.\par\medskip
\noindent\textit{Editorial note.} The counted statement is carried by the author's own lemma environment in the chapter source.\par\medskip
\noindent\textit{Editorial note.} Pre-open gate: python operations/agent-authoring/tools/precard\_gate.py --source weil.tex --label lemma-otimes --start 3752 --end 3784 --local-id 28156 returned RESULT=OK (L28156 free, occupancy FREE, public interval clear). Pre-assembly tag self-check: canonical tag resolved and not present in reference/public-items.jsonl.\par\medskip
\noindent\textit{Editorial note.} The statement and proof are copied from the pinned chapter file; the cited inputs are kept as external pointers with their canonical Stacks tags.\par\medskip
\noindent\textit{Allowed standard prerequisites (editorial).}\begin{itemize}
\item The cohomology of a zero dimensional scheme and the structure of its degree zero part. Same chapter, the lemma describing the cohomology of a smooth projective scheme of dimension zero, including the dimension count and the identification of the degree zero part with a separable algebra. Identifies the degree zero parts appearing in the tensor decomposition with products of coefficient fields.
\item Additivity of the theory under disjoint unions. Same chapter, the lemma computing the cohomology of a disjoint union, together with axiom (B)(a) on the behaviour under products. Splits the tensor product of the two degree zero parts according to the decomposition of the product of the spectra.
\end{itemize}\medskip
\ifdefined\hypertarget\hypertarget{sourcepointer-1}{}\fi
\noindent\textit{Source reference.} Original reference 0FHQ: sources/stacks-weil/pinned-original-source.tex, source lines 3432--3432. Prerequisite pointer invoked at proof line(s) 3764; it names a standard fact of the same chapter that the counted argument uses, and every step of the counted proof built on it is retained verbatim in the delivered body.\par\medskip
\ifdefined\hypertarget\hypertarget{sourcepointer-2}{}\fi
\noindent\textit{Source reference.} Original reference 0FHJ: sources/stacks-weil/pinned-original-source.tex, source lines 2932--2932. Prerequisite pointer invoked at proof line(s) 3772; it names a standard fact of the same chapter that the counted argument uses, and every step of the counted proof built on it is retained verbatim in the delivered body.\par\medskip
\begin{lemma}
\label{lemma-otimes}
Assume given (D0), (D1), (D2), and (D3) satisfying (A), (B), and (C). Let
$k''/k'/k$ be finite separable algebras and let $X$ be a
smooth projective scheme over $k'$. Then
$$
H^*(X) \otimes_{H^0(\Spec(k'))} H^0(\Spec(k'')) =
H^*(X \times_{\Spec(k')} \Spec(k''))
$$
\end{lemma}

\begin{proof}
We will use the results of Lemma \ref{lemma-dim-0} without further mention.
Write
$$
k' \otimes_k k'' = k'' \times l
$$
for some finite separable $k'$-algebra $l$. Write
$F' = H^0(\Spec(k'))$, $F'' = H^0(\Spec(k''))$, and $G = H^0(\Spec(l))$.
Since $\Spec(k') \times \Spec(k'') = \Spec(k'') \amalg \Spec(l)$ we
deduce from axiom (B)(a) and Lemma \ref{lemma-weil-additive}
that we have
$$
F' \otimes_F F'' = F'' \times G
$$
The map from left to right identifies $F''$ with $F' \otimes_{F'} F''$.
By the same token we have
$$
H^*(X) \otimes_F F'' = H^*(X \times_{\Spec(k')} \Spec(k''))
\times H^*(X \times_{\Spec(k')} \Spec(l))
$$
as modules over $F' \otimes_F F'' = F'' \times G$. This proves the lemma.
\end{proof}

\end{document}
